% Four of the routing algorithms of Routing.cpp and the ordering each one declares
% (the full list is the table tab:noc-routing in chapter 9).
\begin{tikzpicture}[font=\sffamily\scriptsize, >=Latex, x=1mm, y=1mm,
  sr/.style={draw=octpurple, fill=octio, semithick, minimum size=4.2mm, inner sep=0, font=\sffamily\tiny},
  srs/.style={sr, fill=octink, text=white, font=\sffamily\tiny\bfseries},
  pth/.style={->, very thick, octink},
  alt/.style={->, very thick, octamber, densely dashed},
  esc/.style={->, very thick, octteal, densely dotted},
  ttl/.style={font=\sffamily\small\bfseries, text=octink, anchor=west},
  lbl/.style={text=octquiet, align=left, anchor=north west, text width=50mm},
  ord/.style={font=\sffamily\scriptsize\bfseries, anchor=west},
  meshgrid/.pic={
    \foreach \yy in {0,1,2,3} \foreach \xx in {0,1,2} \draw[octquiet] (\xx*10,-\yy*10) -- (\xx*10+10,-\yy*10);
    \foreach \yy in {0,1,2} \foreach \xx in {0,1,2,3} \draw[octquiet] (\xx*10,-\yy*10) -- (\xx*10,-\yy*10-10);
    \foreach \yy in {0,1,2,3} \foreach \xx in {0,1,2,3} {\pgfmathtruncatemacro{\rr}{\yy*4+\xx}\node[sr] at (\xx*10,-\yy*10) {\rr};}
  }]
  % ================= XY =================
  \begin{scope}[shift={(0,0)}]
    \node[ttl] at (-3,9) {\code{XY}};
    \pic {meshgrid};
    \draw[pth] (2.4,-30) -- (7.6,-30); \draw[pth] (12.4,-30) -- (17.6,-30); \draw[pth] (22.4,-30) -- (27.6,-30);
    \draw[pth] (30,-27.6) -- (30,-22.4); \draw[pth] (30,-17.6) -- (30,-12.4); \draw[pth] (30,-7.6) -- (30,-2.4);
    \node[srs] at (0,-30) {12}; \node[srs] at (30,0) {3};
    \node[lbl] at (-3,-39.5) {X to the destination column, then Y (then Z on a 3-D lattice). Deterministic, deadlock-free by construction. Any lattice.};
    \node[ord, text=octteal] at (-3,-36.5) {ordering: pair per class};
  \end{scope}
  % ================= WestFirst =================
  \begin{scope}[shift={(55,0)}]
    \node[ttl] at (-3,9) {\code{WestFirst}};
    \pic {meshgrid};
    % 3 -> 12: west first, all the way
    \draw[pth] (27.6,0) -- (22.4,0); \draw[pth] (17.6,0) -- (12.4,0); \draw[pth] (7.6,0) -- (2.4,0);
    \draw[pth] (0,-2.4) -- (0,-7.6); \draw[pth] (0,-12.4) -- (0,-17.6); \draw[pth] (0,-22.4) -- (0,-27.6);
    % 12 -> 3: adaptive between east and north
    \draw[alt] (2.4,-29) -- (7.6,-29); \draw[alt] (10,-27.6) -- (10,-22.4); \draw[alt] (12.4,-20) -- (17.6,-20);
    \draw[alt] (20,-17.6) -- (20,-12.4); \draw[alt] (22.4,-10) -- (27.6,-10); \draw[alt] (30,-7.6) -- (30,-2.4);
    \node[srs] at (0,-30) {12}; \node[srs] at (30,0) {3};
    \node[lbl] at (-3,-39.5) {Black: 3 to 12 turns west first. Amber: 12 to 3 picks, at each router, the minimal output (E or N) with the most credits. \code{NorthLast}, \code{NegativeFirst} and \code{OddEven} work alike with other turn rules. \code{grid2d} only.};
    \node[ord, text=octred] at (-3,-36.5) {ordering: none (adaptive)};
  \end{scope}
  % ================= AdaptiveMinimal =================
  \begin{scope}[shift={(110,0)}]
    \node[ttl] at (-3,9) {\code{AdaptiveMinimal}};
    \fill[octamber!15] (-3,-33) rectangle (33,3);
    \pic {meshgrid};
    \draw[alt] (2.4,-30) -- (7.6,-30); \draw[alt] (10,-27.6) -- (10,-22.4); \draw[alt] (10,-17.6) -- (10,-12.4);
    \draw[alt] (12.4,-10) -- (17.6,-10);
    \draw[esc] (22.4,-10) -- (27.6,-10); \draw[esc] (30,-7.6) -- (30,-2.4);
    \node[srs] at (0,-30) {12}; \node[srs] at (30,0) {3};
    \node[lbl] at (-3,-39.5) {Amber: any minimal output, most credits on the adaptive VCs. Teal: the adaptive VC was not granted, so the escape route, \code{XY} on \code{escape\_vc}, kept to the end. Needs \code{vcs\_per\_class} $\ge 2$, $\ge 3$ with wraps. Any lattice.};
    \node[ord, text=octred] at (-3,-36.5) {ordering: none (adaptive)};
  \end{scope}
  % ================= Table =================
  \begin{scope}[shift={(0,-76)}]
    \node[ttl] at (-3,9) {\code{Table}};
    \node[sr, minimum size=5mm] (a) at (0,0) {A}; \node[sr, minimum size=5mm] (b) at (15,0) {B};
    \node[sr, minimum size=5mm] (c) at (0,-14) {C}; \node[sr, minimum size=5mm] (d) at (15,-14) {D};
    \node[sr, minimum size=5mm] (e) at (30,-7) {E};
    \foreach \u/\v in {a/b, a/c, b/d, c/d, b/e, d/e} \draw[octquiet] (\u) -- (\v);
    \draw[pth] (c) -- (d); \draw[pth] (d) -- (e);
    \node[lbl] at (-3,-24) {Shortest paths by hops, or by \code{link\_weight}. Ties go to the lowest port index. \code{routing\_via} (\code{dest,neighbour}) and \code{routing\_table} (\code{dest,port}) override cells. Any topology, refused if its channel dependency graph has a cycle.};
    \node[ord, text=octteal] at (-3,-21) {ordering: pair per class};
  \end{scope}
  % ================= torus dateline =================
  \begin{scope}[shift={(55,-76)}]
    \node[ttl] at (-3,9) {\code{XY} with wraps};
    \foreach \xx in {0,1,2,3} \node[sr, minimum size=5mm] (t\xx) at (\xx*10,-4) {\xx};
    \foreach \xx/\nx in {0/1,1/2,2/3} \draw[octquiet] (t\xx) -- (t\nx);
    \draw[octquiet] (t3.north) to[out=60, in=120, looseness=0.5] (t0.north);
    \node[text=octquiet] at (15,5.5) {wrap link};
    \draw[pth] (t1) -- (t2);
    \node[text=octink, anchor=west] at (-3,-13.5) {1 to 2: VC 0};
    \draw[alt] (t3.south) to[out=-60, in=-120, looseness=0.45] (t0.south);
    \node[text=octamber, anchor=west] at (17,-13.5) {3 to 0 wraps: VC 1};
    \node[lbl] at (-3,-24) {Torus or ring: the shorter way round each wrapped dimension. A leg that crosses the wrap link uses VC 1, others VC 0 (the dateline rule), chosen from source and destination. Needs \code{vcs\_per\_class} $\ge 2$.};
    \node[ord, text=octteal] at (-3,-21) {ordering: pair per class};
  \end{scope}
  % ================= summary =================
  \begin{scope}[shift={(110,-76)}]
    \node[ttl] at (-3,9) {What a protocol can rely on};
    \node[lbl, text width=48mm, text=octink] at (-3,4) {\code{requires\_ordering(s),}\\\code{pair\_per\_class}, set by the directory presets, accepts only an algorithm that keeps two packets of one class from one source to one destination in order:\\[0.6ex]
      \code{Table} or \code{XY}, \code{vcs\_per\_class} 1\\
      \code{XY} on a torus or a ring (the dateline fixes the VC)\\[0.6ex]
      Anything else is refused at start-up, including the turn models, \code{O1TURN}, \code{Valiant} and \code{ROMM}. No algorithm orders packets of different classes.};
  \end{scope}
\end{tikzpicture}
